% Derivation body (session P0-W3.1). SI-ready: \input-able into the SI or the paper.
% Requirements of the including document: amsmath, array, and \input of theory/macros.tex.
% Labels carry the prefix th: so they do not collide with the paper's labels.
% Symbolic only: no literature value appears in this file. Numbers reach paper text only
% through generated macros.

\section{Regimes and what each link must achieve}\label{th:sec:regimes}

\textit{Physics.} The proposed mechanism is assembled in one chemical environment and must
survive in another. The template is anchored, binds Ca$^{2+}$ and templates the mineral while
the electrode is a water-processed, alkaline slurry at room temperature (\Rs). The finished
interface then sits in a non-aqueous, fluorine-free electrolyte with the Al collector held at
the cathode's anodic potential (\Rc). The two regimes have different solvents, different
reference states and different corrosion modes, so a relation is meaningful only in the regime
whose reference states it uses. Every relation below is therefore tagged with its regime, and
none is carried across without saying so.

Table~\ref{th:tab:regimes} states, for each link, the necessary condition the link must meet
in its regime. These are directions of inequalities, not thresholds: every numeric margin is
registered at G2. Meeting them makes the mechanism thermodynamically possible under the
stated framing; it does not show that a barrier works, and no rate, capacity or device statement
follows from any of them.

\begin{table}[h]
\centering\small
\begin{tabular}{@{}l l >{\raggedright\arraybackslash}p{\dimexpr\textwidth*33/99\relax} >{\raggedright\arraybackslash}p{\dimexpr\textwidth*33/99\relax}@{}}
\hline
Link & Regime & What the link must achieve & Relation\\
\hline
C1 anchor & \Rs & template held at the oxide against water & $\Wwet>0$, Eq.~\eqref{th:eq:dupre}\\
C2 sequester & \Rs & Ca$^{2+}$ enriched at the layer against hydrated Ca$^{2+}$ & $\DGseq<0$, $\Eenr>1$, Eq.~\eqref{th:eq:enr}\\
C3 nucleate & \Rs & heterogeneous barrier lowered, and accessible at $\Oml$ & $f(\theta)$ small; $\DGhet(\Oml)$ low, Eq.~\eqref{th:eq:het}\\
C4 conformal & \Rs, \Rc & a continuous film that closes below the thickness bound & $\Ssp\ge0$ or closure below $d_\mathrm{max}$, Eqs.~\eqref{th:eq:spread}, \eqref{th:eq:combined}\\
C5 stable & \Rs & mineral and oxide on the aqueous Pourbaix hull & $\DGpbx=0$ (tolerance at G2), Eq.~\eqref{th:eq:pbx}\\
C5 stable & \Rc & phases inside their Li-referenced window & $E_\mathrm{red}<\ELi<E_\mathrm{ox}$, Eq.~\eqref{th:eq:window}\\
\hline
\end{tabular}
\caption{Necessary conditions per link. C1--C3 are defined only in \Rs: their reference states
(water-wetted surfaces, hydrated Ca$^{2+}$, an aqueous ion activity product) do not exist in the
cell.}\label{th:tab:regimes}
\end{table}

Two rules follow from the regime split. First, C1--C3 make no statement about the cell:
anchor retention, Ca$^{2+}$ binding and nucleation in the organic electrolyte are not computed.
Second, C5 needs two different constructions (Section~\ref{th:sec:stability}), because aqueous
Pourbaix analysis assumes water as the solvent and the hydrogen electrode as the reference.

\section{Adhesion in water}\label{th:sec:adhesion}

\textit{Physics.} In the slurry the template never meets a dry oxide. The hydroxylated
Al$_2$O$_3$/AlOOH surface is covered by hydrogen-bonded water, and the polymer is hydrated too.
A catechol--oxide contact can only form by displacing that water, so the relevant comparison is
between two states that actually compete in the slurry: the oxide and polymer in contact, and
the oxide and polymer each wetted by water. The dry (vacuum) work of adhesion compares the
contact with two bare surfaces, a state that never occurs in \Rs, and is therefore only a
comparator.

The reversible work per unit area to separate a contact between phases $A$ and $B$ inside a
medium $M$ is the Dupr\'e work $W=\gamma_{AM}+\gamma_{BM}-\gamma_{AB}$: the two new $A$--$M$ and
$B$--$M$ interfaces are created and the $A$--$B$ interface is destroyed. With $A$ the oxide,
$B$ the polymer and $M$ water,
\begin{equation}\label{th:eq:dupre}
\Wwet=\gow+\gpw-\gop \un{J\,m^{-2}}.
\end{equation}
Units: each term is an interfacial free energy, so $\Wwet$ is in J\,m$^{-2}$. Writing every
in-water energy through the dry works of each surface with water,
$\gow=\go+\gw-W_\mathrm{ow}$ and $\gpw=\gp+\gw-W_\mathrm{pw}$, gives
\begin{equation}\label{th:eq:dupre-dry}
\Wwet=\Wdry+2\gw-W_\mathrm{ow}-W_\mathrm{pw},\qquad \Wdry=\go+\gp-\gop \un{J\,m^{-2}}.
\end{equation}
Equation~\eqref{th:eq:dupre-dry} is the water-competition statement in one line: a large dry
adhesion is offset by the works $W_\mathrm{ow}$ and $W_\mathrm{pw}$ with which water already
binds each surface. A hydrophilic, hydroxylated oxide and a hydrophilic polymer bind water
strongly, so $\Wwet$ can be small or negative even when $\Wdry$ is large. The bare-oxide control
has $\Wwet=0$ by construction and fixes the zero of the scale.

The interfacial energies of Eq.~\eqref{th:eq:dupre} are not computed separately. The potential
of mean force (PMF) of a template fragment along a separation coordinate in explicit water gives
the free-energy difference between the separated and bound states, $\Delta F_\mathrm{PMF}$
(J per fragment), which already contains the displacement penalty and the entropy of the
released water. Normalising by the contact area,
\begin{equation}\label{th:eq:pmf}
\Wwet=\frac{\Delta F_\mathrm{PMF}}{A_\mathrm{c}} \un{J\,m^{-2}},
\end{equation}
with units J/m$^2$. Two conditions make Eq.~\eqref{th:eq:pmf} meaningful: the PMF must reach a
plateau at large separation (both surfaces fully rewetted), and $A_\mathrm{c}$ must be defined
once for all fragments, because the ranking across binders depends on it (open question Q5).
C1 requires $\Wwet>0$ beyond its uncertainty on the populated oxide terminations.

\section{The sequestration cycle and the enrichment factor}\label{th:sec:sequestration}

\textit{Physics.} The template can raise the local Ca$^{2+}$ concentration only if it holds
Ca$^{2+}$ more strongly than water does. The question is a competition between two ligand
environments for the same ion: the hydration shell in bulk solution and the coordinating groups
of the template. The free energy must therefore start from hydrated Ca$^{2+}$ at the regime
activity, include the cost of preparing the binding site at the regime pH, and end in a
well-defined local state.

\paragraph{Cycle and standard states.} Let the coordinating group be an acid L--H that binds
Ca$^{2+}$ only when deprotonated. The cycle is
\begin{align}
\text{(i)}\quad &\text{L--H (layer)} \;\to\; \text{L}^{-}\text{(layer)} + \text{H}^{+}\text{(aq)},
&&\Delta G_\mathrm{i}=-\Rg T\ln\alpha \un{J\,mol^{-1}}\nonumber\\
\text{(ii)}\quad &\text{Ca}^{2+}\text{(aq, bulk)} + \text{L}^{-} \;\to\; \text{Ca}^{2+}\!\cdot\text{L}^{-}\text{(layer)},
&&\Delta G_\mathrm{ii}=\DGsite \un{J\,mol^{-1}}\nonumber
\end{align}
with $\alpha=1/\big(1+10^{\,\mathrm{p}K_\mathrm{a}-\mathrm{pH}}\big)$ (dimensionless).
Step (i) is the free energy of finding the site deprotonated when the solution fixes the proton
activity: $\alpha$ is the deprotonated fraction at the regime pH, and $-\Rg T\ln\alpha\ge0$. Step
(ii) is the transfer of one Ca$^{2+}$ from bulk solution into the binding layer. Summing,
\begin{equation}\label{th:eq:cycle}
\DGseq=\DGsite+\Rg T\ln\!\big(1+10^{\,\mathrm{p}K_\mathrm{a}-\mathrm{pH}}\big) \un{J\,mol^{-1}}.
\end{equation}
Units: $\Rg T$ is in J\,mol$^{-1}$ and the logarithm is of a pure number. The correction vanishes
as $\mathrm{pH}-\mathrm{p}K_\mathrm{a}\to\infty$, equals $\Rg T\ln2$ at $\mathrm{pH}=\mathrm{p}K_\mathrm{a}$, and
grows as the site becomes protonated. Both $\mathrm{p}K_\mathrm{a}$ and the activity coefficients are
taken at the regime ionic strength. This is why catechol, largely protonated near neutral pH,
must be tested rather than assumed (design flag 4).

Step (ii) fixes the standard states. Let $w(\mathbf r)$ be the PMF of one Ca$^{2+}$ relative to
bulk solution (J per ion, $w\to0$ far from the template), computed with counter-ions and water
explicit. In the dilute limit the equilibrium concentration is $c(\mathbf r)=\cb\,\e^{-w(\mathbf r)/\kB T}$.
Assigning to each site a volume $\Vsite$ (the binding-layer volume per site) and defining
\begin{equation}\label{th:eq:site}
\DGsite\equiv-\Rg T\ln\!\Big[\frac{1}{\Vsite}\int_{\mathrm{site}}\e^{-w(\mathbf r)/\kB T}\,\dd^3 r\Big] \un{J\,mol^{-1}}
\end{equation}
(the bracket is a ratio of volumes, so it is a pure number) makes the transfer free energy a
partition between two concentration scales: the bulk at the regime activity on the solution
side, and the mean concentration in the site volume on the layer side. The conventional binding
free energy to a 1\,mol-per-volume standard state $c^{\circ}$ differs from it by the
standard-state correction
\begin{equation}\label{th:eq:stdstate}
\Delta G^{\circ}_\mathrm{bind}=\DGsite-\Rg T\ln\!\big(\NA c^{\circ}\Vsite\big) \un{J\,mol^{-1}},
\end{equation}
where $\NA c^{\circ}\Vsite$ is dimensionless. A literature binding constant can therefore be
compared with $\DGsite$ only after Eq.~\eqref{th:eq:stdstate}, and only with the same $\Vsite$.
The choice of $\Vsite$ shifts $\DGsite$ by $\Rg T\ln(\Vsite'/\Vsite)$, so it is fixed once
(open question Q4).

\paragraph{Activity versus concentration.} Equation~\eqref{th:eq:site} defines a concentration
ratio. If $w$ contains every interaction of the ion with its surroundings, the chemical
potential of Ca$^{2+}$ is uniform at equilibrium and the activity in the layer equals the bulk
activity; the whole enrichment then sits in a local activity coefficient. The screen instead
uses a \emph{local-ideal} bookkeeping: the local Ca$^{2+}$ activity is the local concentration
times the bulk activity coefficient, and the binding free energy is carried explicitly by
$\DGseq$. This is a definition, not an approximation, provided the binding free energy is not
counted a second time elsewhere (Section~\ref{th:sec:template}).

\paragraph{Enrichment and its validity range.} With the local-ideal bookkeeping and dilute,
non-interacting sites,
\begin{equation}\label{th:eq:enr}
\Eenr\equiv\frac{\cl}{\cb}=\exp\!\Big(-\frac{\DGseq}{\Rg T}\Big) \un{1}.
\end{equation}
$\DGseq<0$ gives enrichment ($\Eenr>1$); $\DGseq>0$ gives depletion: the template is then a worse
ligand than water. Equation~\eqref{th:eq:enr} holds when (a) Ca$^{2+}$ is dilute in bulk, (b)
the sites are far from saturation, (c) one ion binds per site and sites do not interact, and (d)
the local and bulk activity coefficients are equal. Condition (b) is quantified next.

\paragraph{Saturation bound.} The layer has $\Gs$ sites per unit area in a thickness $\ds$, so
its Ca$^{2+}$ content cannot exceed $\cmax=\Gs/\ds$ (mol\,m$^{-2}$ over m gives mol\,m$^{-3}$).
For independent, equivalent sites (Langmuir), the occupied fraction is
$\Theta=x/(1+x)$ with $x=\Eenr\cb/\cmax$, and $\cl=\cmax\Theta$:
\begin{equation}\label{th:eq:langmuir}
\cl=\frac{\Eenr\cb}{1+\Eenr\cb/\cmax}\le\min\big(\Eenr\cb,\ \cmax\big) \un{mol\,m^{-3}}.
\end{equation}
The dilute result \eqref{th:eq:enr} is the limit $x\ll1$; at $x\gg1$ the enrichment saturates at
$\cmax/\cb$ however negative $\DGseq$ is. A strongly binding but sparse template is therefore
bounded by its site density, not by its binding free energy. Electrostatic coupling between
charged sites, which Langmuir ignores, is open question Q3.

\section{Local supersaturation}\label{th:sec:supersaturation}

\textit{Physics.} Crystallisation is driven by the excess chemical potential of the dissolved
ions over the solid. For a 1:1 salt that excess is set by the product of the two ion activities
relative to its value at saturation. The template can change that product only where the
nucleus forms, and only through the ion it binds.

One definition is used throughout:
\begin{equation}\label{th:eq:omega}
\Om=\frac{\mathrm{IAP}}{\Ksp}=\frac{\aCa\,\aCO}{\Ksp} \un{1},\qquad \Dmu=\kB T\ln\Om \un{J},
\end{equation}
with $\Ksp$ the solubility product of the registered framing (ACC or calcite) and $\Dmu$ the
chemical-potential drop per formula unit on crystallisation. Units: activities and $\Ksp$ are
pure numbers on the molal scale; $\kB T$ is in J. Nucleation requires $\Om>1$. The per-ion form
$S=\Om^{1/2}$, $\Dmu=2\kB T\ln S$, is the same statement and is not used, so that no factor of
two can be lost between lanes. The charter's $S$ for C3 is this $\Om$.

\paragraph{Enrichment of Ca$^{2+}$ alone.} In the local-ideal bookkeeping the local Ca$^{2+}$
activity is $\aCa\,\cl/\cb$, while the carbonate activity is held at its bulk value (carbonate is
not assumed enriched). Then
\begin{equation}\label{th:eq:omlocal}
\Oml=\frac{\cl}{\cb}\,\Omb\;\xrightarrow{\ \Eenr\cb\ll\cmax\ }\;\Eenr\,\Omb \un{1},\qquad
\Dmu_\mathrm{local}=\Dmu_\mathrm{bulk}-\frac{\DGseq}{\NA} \un{J}.
\end{equation}
Units: $\DGseq/\NA$ is J per ion, matching $\Dmu$ per formula unit, since one Ca$^{2+}$ enters
each formula unit. The saturation bound carries over: $\Oml\le(\cmax/\cb)\,\Omb$.

Two cautions belong to this step. (1) A template that attracts Ca$^{2+}$ is usually anionic and
repels CO$_3^{2-}$. Holding $\aCO$ at its bulk value is then an upper estimate of $\Oml$, biased in
favour of the mechanism; the bound is acceptable for a screen only if it is reported as such (Q2).
(2) Because the true chemical potential is uniform at equilibrium, the increase of $\Om$ from
$\Omb$ to $\Oml$ is the binding free energy moved into the driving force by the bookkeeping. It
is legitimate only if the same free energy is not also credited through the interfacial
energies (Section~\ref{th:sec:template}).

\section{Homogeneous and heterogeneous nucleation}\label{th:sec:nucleation}

\textit{Physics.} A cluster of the new phase gains $\Dmu$ for every formula unit it contains but
pays $\gamma$ for every unit of interface it creates. Volume grows as $r^{3}$ and surface as
$r^{2}$, so small clusters are unstable and large ones grow; the barrier is the free energy at the
crossover. On a substrate the cluster can replace part of its interface with a cheaper
substrate contact, which lowers the barrier without changing the driving force.

\paragraph{Homogeneous barrier.} In the capillarity approximation, a spherical cluster of radius
$r$ with molecular volume $\vm=\Vm/\NA$ has
\begin{equation}\label{th:eq:dg-r}
\Delta G_\mathrm{hom}(r)=-\frac{4\pi r^{3}}{3}\frac{\Dmu}{\vm}+4\pi r^{2}\gnl \un{J}.
\end{equation}
Units: $r^{3}/\vm$ counts formula units, times J per formula unit; $r^{2}\gnl$ is m$^2$ times
J\,m$^{-2}$. Setting $\partial\Delta G_\mathrm{hom}/\partial r=0$,
\begin{equation}\label{th:eq:hom}
r^{*}=\frac{2\gnl\vm}{\Dmu} \un{m},\qquad
\DGhom=\frac{16\pi\gnl^{3}\vm^{2}}{3\Dmu^{2}}=\frac{4\pi r^{*2}\gnl}{3} \un{J}.
\end{equation}
Units: J\,m$^{-2}\cdot$m$^{3}$/J $=$ m for $r^{*}$; J$^{3}$\,m$^{-6}\cdot$m$^{6}$/J$^{2}$ $=$ J for
$\DGhom$. With $\Dmu=\kB T\ln\Om$, $\DGhom\to\infty$ as $\Om\to1^{+}$ (no driving force, no
nucleation) and $\DGhom\to0$ as $\Om\to\infty$; it increases with $\gnl$ and $\Vm$ and decreases
with $T$ at fixed $\Om$. The absolute barrier in units of $\kB T$ is $\DGhom/\kB T$.

\paragraph{Spherical-cap geometry.} On a flat substrate the cluster is a cap cut from a sphere of
radius $r$, meeting the substrate at contact angle $\theta$ measured through the cap. Elementary
geometry gives base radius $\rho=r\sin\theta$, height $h=r(1-\cos\theta)$, and
\begin{align}
A_\mathrm{cap}&=2\pi r^{2}(1-\cos\theta),\qquad A_\mathrm{base}=\pi r^{2}\sin^{2}\theta \un{m^{2}},\label{th:eq:cap-areas}\\
V_\mathrm{cap}&=\frac{\pi r^{3}}{3}(1-\cos\theta)^{2}(2+\cos\theta) \un{m^{3}}.\label{th:eq:cap-volume}
\end{align}
The cap's free energy has three parts: the volume term, the free nucleus--liquid surface, and the
base, where substrate--liquid interface ($\gsl$) is replaced by substrate--nucleus interface
($\gsn$):
\begin{equation}\label{th:eq:dg-cap}
\Delta G_\mathrm{het}(r)=-V_\mathrm{cap}\frac{\Dmu}{\vm}+\gnl A_\mathrm{cap}+(\gsn-\gsl)A_\mathrm{base} \un{J}.
\end{equation}
Stationarity of Eq.~\eqref{th:eq:dg-cap} with respect to $\theta$ at fixed cap volume is the
contact-line balance, Young's relation $\gsl-\gsn=\gnl\cos\theta$ (Section~\ref{th:sec:young}).
Substituting it, the surface terms become
\begin{equation}
\gnl\pi r^{2}\big[2(1-\cos\theta)-\cos\theta\sin^{2}\theta\big]
=\gnl\pi r^{2}(1-\cos\theta)^{2}(2+\cos\theta)=4\pi r^{2}\gnl\,f(\theta) \un{J},
\end{equation}
using $2-\cos\theta-\cos^{2}\theta=(1-\cos\theta)(2+\cos\theta)$, and the volume term is
$\frac{4\pi r^{3}}{3}f(\theta)$ times $-\Dmu/\vm$, with
\begin{equation}\label{th:eq:f}
f(\theta)=\frac{(2+\cos\theta)(1-\cos\theta)^{2}}{4}=\frac{2-3\cos\theta+\cos^{3}\theta}{4} \un{1}.
\end{equation}
Hence $\Delta G_\mathrm{het}(r)=f(\theta)\,\Delta G_\mathrm{hom}(r)$ at every $r$: the cap has the
same critical radius as the sphere, and
\begin{equation}\label{th:eq:het}
\DGhet=f(\theta)\,\DGhom \un{J}.
\end{equation}
The factorisation holds only because one and the same $\gnl$ multiplies $A_\mathrm{cap}$ and enters
Young's relation; this is the reason Section~\ref{th:sec:young} keeps the bare $\gnl$ in Young.

\paragraph{Properties of $f$.} $f(0)=0$ (complete wetting, no barrier), $f(\pi/2)=\tfrac12$
(hemisphere), $f(\pi)=1$ (no wetting, the homogeneous barrier). Its derivative is
$f'(\theta)=\tfrac34\sin^{3}\theta\ge0$ on $[0,\pi]$, so $f$ is monotone and
$0\le\DGhet\le\DGhom$ for every admissible $\theta$. $f$ is also the volume fraction of the cap
relative to the full sphere of the same radius.

\paragraph{Assumptions.} Capillarity (macroscopic interfacial energies at nucleus size), an
isotropic $\gnl$ (closer to the amorphous framing than to faceted calcite), a flat, homogeneous
substrate, no line tension, and a cluster of the bulk phase's composition. A non-classical
pathway (prenucleation clusters, liquid precursors) is outside this framework and is a
limitation, not a parameter.

\section{Young's relation and the fixed sign convention}\label{th:sec:young}

\textit{Physics.} At the contact line, moving the line outward by a small amount replaces
substrate--liquid interface by substrate--nucleus interface and stretches the nucleus--liquid
surface. At equilibrium these balance. The contact angle therefore measures how much the
substrate prefers the nucleus to the liquid, in units of the nucleus--liquid energy.

\paragraph{Convention, in words.} $\theta$ is the angle at the three-phase line measured through
the nucleus. $\cos\theta$ is the free energy gained per unit contact area by replacing
substrate--liquid interface with substrate--nucleus interface, divided by the nucleus--liquid
interfacial free energy. $\theta=0$ means the nucleus wets the template completely; $\theta=\pi$
means it does not wet it at all.

\paragraph{Convention, as an equation.}
\begin{equation}\label{th:eq:young}
\cos\theta=\frac{\gsl-\gsn}{\gnl},\qquad \gsn=\gsnz+\dgw \un{1;\ J\,m^{-2}}.
\end{equation}
Units: a ratio of interfacial free energies. A finite angle exists only for
$|\gsl-\gsn|\le\gnl$. For $\gsl-\gsn\ge\gnl$ the nucleus phase spreads ($\Ssp\ge0$,
Section~\ref{th:sec:conformality}) and $f=0$; for $\gsl-\gsn\le-\gnl$ it does not wet and $f=1$.
An implementation clips $\cos\theta$ to $[-1,1]$ and reports which branch applied.

\paragraph{The water-displacement term.} The template is hydrated. When the nucleus forms on it,
the template-bound water layer under the contact disc must be displaced, and that costs free
energy over the contact area only. Define
\begin{equation}\label{th:eq:dgw}
\dgw\equiv\gsn-\gsnz \un{J\,m^{-2}},
\end{equation}
where $\gsnz$ is a substrate--nucleus energy from a reference that does not include the removal
of that layer (for example, a contact computed without the template's structured water). The
sign convention is: \textbf{$\dgw>0$ is a cost.} From Eq.~\eqref{th:eq:young},
$\partial\cos\theta/\partial\dgw=-1/\gnl<0$, and with $\partial f/\partial\cos\theta=-\tfrac34\sin^{2}\theta\le0$,
\begin{equation}\label{th:eq:signflip}
\frac{\partial f}{\partial\dgw}=\frac{3\sin^{2}\theta}{4\gnl}\ge0 \un{m^{2}\,J^{-1}}:
\end{equation}
a positive displacement cost raises $\theta$, $f$ and $\DGhet$, and $\dgw<0$ lowers them. If
$\gsn$ is supplied directly as a full in-water interfacial free energy, the displacement is
already inside it, $\dgw\equiv0$, and the term must not be added again. Applying it twice is the
double-counting error that C3 item 13 demonstrates.

\paragraph{Why Young uses the bare $\gnl$ and not $\geff$.} The deck writes
$\geff=\gnl-\dgw$ and places it in the denominator of Young. That is rejected, for three reasons.
(1) Geometry: the displaced water sits on the contact disc. The free surface of the cap, whose
energy is $\gnl A_\mathrm{cap}$, never touches the template, so nothing about the template's
hydration changes $\gnl$. (2) Consistency: the factorisation of Eq.~\eqref{th:eq:het} requires the
$\gnl$ in Young to equal the $\gnl$ that multiplies $A_\mathrm{cap}$ and sets $\DGhom$. A modified
denominator would leave a residual surface term and break $\DGhet=f\DGhom$. (3) Sign: with
$\geff$ in the denominator,
\begin{equation}\label{th:eq:deck-sign}
\frac{\partial\cos\theta}{\partial\dgw}\bigg|_\mathrm{deck}=\frac{\gsl-\gsnz}{(\gnl-\dgw)^{2}} \un{m^{2}\,J^{-1}},
\end{equation}
which is positive on the wetting side ($\gsl>\gsnz$): a displacement \emph{cost} would lower the
barrier. This is design flag 3. $\geff$ is kept only as a named diagnostic, so that W4 can
reproduce the deck's form for the double-counting demonstration. The contract encodes the
reversal as sign-flip test SF1 (Table~\ref{th:tab:signflip}).

\begin{table}[h]
\centering\small
\begin{tabular}{@{}l l l@{}}
\hline
Change & Registered convention & Deck form ($\geff$ in the denominator), $\gsl>\gsnz$\\
\hline
$\dgw>0$ & $\cos\theta\downarrow$, $\DGhet\ge\DGhet|_{\dgw=0}$ & $\cos\theta\uparrow$, $\DGhet\le\DGhet|_{\dgw=0}$\\
$\dgw<0$ & $\DGhet\le\DGhet|_{\dgw=0}$ & reversed\\
$\DGseq<0$ & $\Oml>\Omb$, barrier lowered & (not affected by the Young form)\\
$\DGseq>0$ & $\Oml<\Omb$, barrier raised & \\
\hline
\end{tabular}
\caption{Sign-flip tests SF1 and SF2 of \texttt{theory/contract.yaml}.}\label{th:tab:signflip}
\end{table}

\section{Where the template acts, and the no-double-counting condition}\label{th:sec:template}

\textit{Physics.} The template is both an interface and a ligand field. As an interface it
changes the energies of the surfaces the nucleus replaces and creates; as a ligand field it
concentrates Ca$^{2+}$ where the nucleus forms. In classical nucleation the barrier depends on
exactly four quantities, $(\gsl,\gsn,\gnl,\Dmu)$, so the template can act only through $\theta$
and through $\Om$. The risk is that one microscopic interaction, the binding of Ca$^{2+}$ to the
template, is counted in both.

\paragraph{Statement.} Combining Eqs.~\eqref{th:eq:omlocal}, \eqref{th:eq:hom}, \eqref{th:eq:het}
and \eqref{th:eq:young},
\begin{equation}\label{th:eq:composite}
\DGhet=f\big(\theta[\gsl,\gsnz,\dgw]\big)\,
\frac{16\pi\gnl^{3}\vm^{2}}{3\big(\kB T\ln\Oml[\DGseq,\Gs,\ds]\big)^{2}} \un{J}.
\end{equation}
\textbf{No-double-counting condition.} (a) Every template-dependent input enters exactly one
channel: $\gsl,\gsnz,\dgw$ only $\theta$; $\DGseq,\Gs,\ds$ only $\Om$; so
$\partial\theta/\partial\DGseq=0$ and $\partial\Om/\partial\gamma_{\mathrm{s}\cdot}=0$. Bulk inputs
($\gnl$, $\Vm$, $T$) may enter both. (b) The interfacial energies $\gsl$ and $\gsn$ are evaluated
for the template \emph{without} the Ca$^{2+}$ excess that $\Oml$ already credits. Part (a) is
checked mechanically by \texttt{theory/tools/check\_contract.py}; part (b) is the physical content
and is proved in P0-W3.2.

\paragraph{Outline of the proof (P0-W3.2).}
\begin{enumerate}
\item \emph{Exact reference.} Write the free energy of a cap on the template in the grand
ensemble at the bulk chemical potentials of Ca$^{2+}$ and CO$_3^{2-}$. In this bookkeeping (A)
the driving force is $\Dmu_\mathrm{bulk}$ everywhere, and all template--ion interactions are
inside grand-canonical interfacial energies $\gsl^{(\mu)}$, $\gsn^{(\mu)}$.
\item \emph{Gibbs adsorption.} At fixed $T$, $\dd\gsl^{(\mu)}=-\sum_i\Gamma_i\,\dd\mu_i$, so the
Ca$^{2+}$ excess $\Gamma_\mathrm{Ca}$ at the template lowers $\gsl^{(\mu)}$ relative to the
Ca-free template by $\int\Gamma_\mathrm{Ca}\,\dd\mu_\mathrm{Ca}$; the same holds at the
substrate--nucleus interface with its own excess.
\item \emph{Local-ideal bookkeeping (B), the one registered here.} Move the binding free energy
out of the interfacial energies and into the driving force: $\Dmu_\mathrm{local}$ from
Eq.~\eqref{th:eq:omlocal}, with $\gsl,\gsn$ of the Ca-free template.
\item \emph{Equivalence or bound.} Show under which conditions (A) and (B) give the same
$\DGhet$, or bound their difference. The key quantity is the Ca$^{2+}$ excess at the
substrate--liquid interface \emph{minus} that at the substrate--nucleus interface: in (A),
adsorption that stabilises only the template--water interface \emph{raises} $\theta$; it lowers
the barrier only if the nucleus incorporates the bound ions at a lower cost than water holds
them.
\item \emph{Error demonstration.} Mixing (A)'s Ca-loaded $\gsl$ with (B)'s $\Oml$ counts the same
binding twice; the size of that error is C3 item 13, together with the $\dgw$ double application
of Section~\ref{th:sec:young}.
\end{enumerate}
Step 4 matters for the screen, and is flagged as an open question (Q1), because bookkeeping (B)
always lowers the barrier when $\DGseq<0$, while in (A) the sign of the template's effect depends
on a difference of adsorption excesses. The choice of bookkeeping is therefore not neutral, and
the proof must settle it before any C3 output is read.

\section{Film conformality}\label{th:sec:conformality}

\textit{Physics.} A barrier must be continuous: one through-path lets electrolyte reach Al. A
low nucleation barrier says where nuclei form easily, not whether they spread into a film. Whether
the mineral spreads or beads on the template is a wetting balance, and whether beads merge before
they grow thick is a competition between how many islands form and how far each spreads.

\paragraph{Spreading coefficient.} The free-energy change per unit area when a thick layer of the
nucleus phase replaces the template--liquid interface is
\begin{equation}\label{th:eq:spread}
\Ssp=\gsl-\gsn-\gnl=\gnl(\cos\theta-1) \un{J\,m^{-2}},
\end{equation}
the second form holding on the partial-wetting branch where Young has a solution. Units: a sum of
interfacial free energies. Since $\cos\theta\le1$, a finite contact angle always has $\Ssp\le0$.
The sign condition is:
$\Ssp\ge0$ means no finite angle exists, $\theta\to0$, $f\to0$, and the phase grows layer by
layer (Frank--van der Merwe); $\Ssp<0$ means $\theta>0$ and growth by islands (Volmer--Weber).
Island-then-layer growth (Stranski--Krastanov) needs a strain or misfit energy that this
isotropic framework does not contain. The registered $\dgw$ enters $\Ssp$ through $\gsn$, with
the same sign as in $\theta$: a displacement cost makes spreading harder.

\paragraph{Coalescence of islands.} With $\Ssp<0$, a film forms only when islands meet. Let the
islands be equal caps of contact angle $\theta$ on a square lattice of density $\Ns$ (spacing
$\ell=\Ns^{-1/2}$). A cap of base radius $\rho$ has height $h=\rho\tan(\theta/2)$. Neighbours touch
when $2\rho=\ell$, and the base discs then cover only $\pi/4$ of the area; the square cells are
covered when $\rho=\ell/\sqrt2$. The two geometric thicknesses are
\begin{equation}\label{th:eq:hc}
\hc=\frac{\tan(\theta/2)}{2\sqrt{\Ns}},\qquad
h_\mathrm{cl}=\frac{\tan(\theta/2)}{\sqrt{2\Ns}}=\sqrt2\,\hc \un{m}.
\end{equation}
Units: $\Ns^{-1/2}$ is a length. Both vanish as $\theta\to0$ and diverge as $\theta\to\pi$. They
are geometric estimates that treat caps as growing self-similarly up to contact; they are
arguments, not measurements, and $\Ns$ comes from a nucleation-and-growth model (item 17), not
from thermodynamics.

\paragraph{What a contact angle does and does not imply about coverage.} A small $\theta$ means
thin islands and early closure for a given $\Ns$. It does not fix $\Ns$, which depends on the
nucleation rate relative to lateral growth; it says nothing about places where the template is
absent, where the bare-oxide angle applies; and it is silent on porosity, on dissolution and
reprecipitation of the amorphous phase, and on cracking when a hydrated film shrinks on drying.
The proposed C3 band $f(\theta)\le\tfrac12$ admits $\theta$ up to $\pi/2$, which is island growth;
conformality is therefore a separate criterion (C4, design flag 2), chosen at G2 among a
small-angle bound, $\Ssp\ge0$, and a closure condition such as Eq.~\eqref{th:eq:combined}.

\section{Thermodynamic stability}\label{th:sec:stability}

\textit{Physics.} A layer that is thermodynamically unstable in its own environment dissolves or
reacts. The environment differs between regimes, so stability is two questions: in the slurry,
whether CaCO$_3$ and the native oxide survive alkaline water at the regime pH and potential; in
the cell, whether the phases survive the Li reservoir at the collector's potential. Both are
screens: necessary, not sufficient, and blind to morphology and kinetics.

\paragraph{\Rs: aqueous Pourbaix.} Each species $\phi$ (solid or dissolved) is brought to a
common reference by exchanging O and H with water, protons and electrons at the regime pH and
potential. Per non-O, non-H atom,
\begin{multline}\label{th:eq:pbx}
\Psi_\phi=\frac{1}{n_\mathrm{M}}\Big[G_\phi-n_\mathrm{O}\mu_{\mathrm{H_2O}}-\Rg T\ln10\,(2n_\mathrm{O}-n_\mathrm{H})\,\mathrm{pH}\\
-(2n_\mathrm{O}-n_\mathrm{H}+Q)\Fc\ESHE\Big] \un{J\,mol^{-1}},\qquad \DGpbx=\Psi_\phi-\Psi_\mathrm{hull}\ge0.
\end{multline}
Units: $G_\phi$ and $\mu_{\mathrm{H_2O}}$ are J\,mol$^{-1}$, $\Rg T\ln10\,\mathrm{pH}$ is J\,mol$^{-1}$,
and $\Fc\ESHE$ is C\,mol$^{-1}\cdot$V $=$ J\,mol$^{-1}$. Species with $2n_\mathrm{O}-n_\mathrm{H}+Q>0$
are favoured by raising $\ESHE$, and species with $2n_\mathrm{O}-n_\mathrm{H}>0$ by raising pH.
$\Psi_\mathrm{hull}$ is the lower envelope over all species and mixtures of the same metal
composition, and $\DGpbx=0$ means stable. \emph{What it can say:} whether the phase is the
equilibrium state at the chosen (pH, $\ESHE$) and dissolved-ion activities. \emph{What it cannot
say:} rates, passivation, local pH at the interface, morphology. Two traps: the amorphous framing
lies above the calcite hull by its offset, so ACC is never on the hull and the tolerance must be
registered (C5 open item); and the dissolved-ion activities chosen for the aqueous species decide
whether solids appear stable at all, so the screen can be satisfied by construction unless it is
shown to fail where it must (C5 item 21, Q8).

\paragraph{\Rc: Li-open grand potential.} In the cell the solvent is not water and the reference
is Li$^+$/Li. Holding the collector at $\ELi$ fixes the Li chemical potential,
\begin{equation}\label{th:eq:muli}
\muLi=\muLiz-\Fc\ELi \un{J\,mol^{-1}},\qquad \Phi_\phi=G_\phi-n_\mathrm{Li}\muLi \un{J\,mol^{-1}}.
\end{equation}
For a reaction $\phi\to\sum_k\nu_k\psi_k+\Delta n\,\mathrm{Li}$ that releases $\Delta n$ Li to the
reservoir, Li conservation gives
$\Delta\Phi(\ELi)=\sum_k\nu_k\Phi_k-\Phi_\phi=\Delta G_\mathrm{r}-\Delta n\,\Fc\ELi$, with
$\Delta G_\mathrm{r}$ the reaction free energy against Li metal. The phase is stable against
that reaction while $\Delta\Phi>0$, so each reaction sets one limit,
\begin{equation}\label{th:eq:window}
E^{*}=\frac{\Delta G_\mathrm{r}}{\Delta n\,\Fc} \un{V},\qquad
\begin{cases}\Delta n>0\ (\text{delithiation, oxidation}): & \ELi<E^{*}\ \to\ E_\mathrm{ox}=\min E^{*},\\
\Delta n<0\ (\text{lithiation, reduction}): & \ELi>E^{*}\ \to\ E_\mathrm{red}=\max E^{*}.\end{cases}
\end{equation}
Units: J\,mol$^{-1}$ over C\,mol$^{-1}$ is V. \emph{What it can say:} whether a phase reacts with
Li at the collector's potential. \emph{What it cannot say:} oxidation that exchanges no Li. A
Li-free phase such as CaCO$_3$ has no reaction with $\Delta n>0$, so its Li-open $E_\mathrm{ox}$ is
unbounded: the construction declares it stable at any anodic potential by construction. Anodic
stability of CaCO$_3$ and of the Al phases needs oxidation reactions that release CO$_2$, O$_2$ or
dissolved Ca$^{2+}$ and Al$^{3+}$ into the electrolyte, which require the electrolyte species as
additional open components (Q9). Kinetics and electrolyte oxidation are outside both screens.

\section{The insulating-film limitation}\label{th:sec:film}

\textit{Physics.} CaCO$_3$ is an electronic insulator. The film lies between the cathode coating and
the Al collector, and every electron of the cathode current must cross it. A film that is
continuous but too thick blocks the collector, so conformality cannot be bought with thickness.
Two transport limits bound the admissible thickness; the numbers come later, through macros.

\paragraph{Ohmic limit.} If electrons cross by bulk conduction with resistivity $\rhoe$, the
area-specific resistance and the bound set by an admissible $R_{A,\max}$ are
\begin{equation}\label{th:eq:ohm}
R_A=\rhoe d \un{\Omega\,m^{2}},\qquad d\le d_\mathrm{max}^\mathrm{ohm}=\frac{R_{A,\max}}{\rhoe} \un{m}.
\end{equation}
Units: $\Omega$\,m times m is $\Omega$\,m$^{2}$; $\Omega$\,m$^{2}$ over $\Omega$\,m is m.

\paragraph{Tunnelling limit.} For a thin film the current is carried by direct tunnelling through
a barrier of height $\phiB$. In the WKB approximation for a rectangular barrier the transmission
decays as $\e^{-2\kappa d}$, so
\begin{equation}\label{th:eq:tun}
j=j_0\,\e^{-2\kappa d} \un{A\,m^{-2}},\qquad \kappa=\frac{\sqrt{2m_\mathrm{e}\phiB}}{\hbar} \un{m^{-1}},\qquad
d\le d_\mathrm{max}^\mathrm{tun}=\frac{\ln(j_0/\jreq)}{2\kappa} \un{m}.
\end{equation}
Units: $\sqrt{\mathrm{kg\cdot J}}/(\mathrm{J\,s})=\mathrm{m}^{-1}$; the logarithm is of a ratio of current
densities. The bound is positive only for $j_0>\jreq$. The prefactor $j_0$, the effective mass and
the barrier height are inputs with large uncertainty (Q10), and defects or hopping can only raise
the current, so Eq.~\eqref{th:eq:tun} is a conservative bound for a defect-free film.

\paragraph{Combined with conformality.} The film must close before it exceeds the admissible
thickness:
\begin{equation}\label{th:eq:combined}
h_\mathrm{cl}=\frac{\tan(\theta/2)}{\sqrt{2\Ns}}\le d_\mathrm{max}=\max\big(d_\mathrm{max}^\mathrm{ohm},\,d_\mathrm{max}^\mathrm{tun}\big) \un{m}.
\end{equation}
This is a stated limitation (C4 item 18), not a performance figure. It contains no statement
about cell resistance, rate or capacity.

\section{Summary: links and the inputs they consume}\label{th:sec:summary}

\textit{Physics.} Each link reduces to a small set of inputs. Listing them once is what lets the
parameter table (W1) fetch every input and the implementation (W4) consume every one, with
nothing typed and nothing missing. Table~\ref{th:tab:inputs} is that list; the ids are those of
\texttt{theory/symbols.yaml}, and the functions are those of \texttt{theory/contract.yaml}.

\begin{table}[h]
\centering\small
\footnotesize\begin{tabular}{@{}>{\raggedright\arraybackslash}p{\dimexpr\textwidth*12/99\relax} l >{\raggedright\arraybackslash}p{\dimexpr\textwidth*22/99\relax} >{\raggedright\arraybackslash}p{\dimexpr\textwidth*29/99\relax} >{\raggedright\arraybackslash}p{\dimexpr\textwidth*16/99\relax}@{}}
\hline
Link & Regime & Contract functions & Inputs consumed (symbol ids) & Source\\
\hline
C1 adhesion & \Rs & \texttt{W\_adh\_wet}, \texttt{W\_adh\_wet\_from\_pmf}, \texttt{W\_adh\_dry\_fn} &
\texttt{Delta\_F\_pmf}, \texttt{A\_c}; comparator \texttt{gamma\_o}, \texttt{gamma\_p}, \texttt{gamma\_op}, \texttt{gamma\_w} &
computed PMF; W1 for comparators\\
C2 sequestration & \Rs & \texttt{Delta\_G\_seq\_cycle}, \texttt{E\_enr}, \texttt{c\_max}, \texttt{c\_local} &
\texttt{Delta\_G\_site}, \texttt{pKa}, \texttt{pH}, \texttt{T}, \texttt{c\_bulk}, \texttt{Gamma\_s}, \texttt{delta\_s}, \texttt{V\_site} &
computed PMF; W1 for pKa, pH, c\_bulk\\
Super\-saturation & \Rs & \texttt{Omega\_bulk}, \texttt{Omega\_local}, \texttt{Delta\_mu} &
\texttt{a\_Ca}, \texttt{a\_CO3}, \texttt{K\_sp}, \texttt{T}, with C2 outputs &
W1 (activities from speciation at the regime pH and ionic strength)\\
C3 nucleation & \Rs & \texttt{r\_star}, \texttt{DG\_hom}, \texttt{f\_theta}, \texttt{DG\_het}, \texttt{cos\_theta\_young}, \texttt{theta\_young}; audit: \texttt{gamma\_eff\_deck}, \texttt{cos\_theta\_deck} &
\texttt{gamma\_nl}, \texttt{gamma\_sl}, \texttt{gamma\_sn0}, \texttt{Delta\_gamma\_w}, \texttt{V\_m}, \texttt{T}, with $\Oml$ &
W1 priors and computed slabs; framing at G2\\
C4 conformality & \Rs & \texttt{S\_spread}, \texttt{h\_c}, \texttt{h\_cl} &
C3 energies, \texttt{N\_s} &
nucleation-and-growth model (item 17)\\
C4 film bound & \Rc & \texttt{R\_A}, \texttt{d\_max\_ohm}, \texttt{kappa}, \texttt{d\_max\_tun} &
\texttt{rho\_e}, \texttt{R\_A\_max}, \texttt{phi\_B}, \texttt{m\_e}, \texttt{j\_0}, \texttt{j\_req}, \texttt{d\_film} &
W1; \texttt{R\_A\_max}, \texttt{j\_req} registered at G2\\
C5 Pourbaix & \Rs & \texttt{Psi\_pbx}, \texttt{Delta\_G\_pbx} &
\texttt{G\_phase} (with the ACC offset), \texttt{mu\_H2O}, \texttt{N\_O}, \texttt{N\_H}, \texttt{N\_M}, \texttt{Q\_charge}, \texttt{pH}, \texttt{E\_SHE}, \texttt{T} &
W1 thermochemistry; materials database\\
C5 window & \Rc & \texttt{mu\_Li}, \texttt{Phi\_Li}, \texttt{E\_window} &
\texttt{G\_phase}, \texttt{N\_Li}, \texttt{mu\_Li0}, \texttt{E\_Li}, \texttt{Delta\_G\_r}, \texttt{Delta\_n\_Li} &
materials database\\
\hline
\end{tabular}
\caption{Each link, the contract functions that implement it, and the inputs it consumes. Every
input is fetched or computed; none is a result of the study.}\label{th:tab:inputs}
\end{table}
